% =========================================================
\chapter{Apêndices}
% =========================================================

O apêndice deve conter material produzido pelo próprio autor que seja
relevante para complementar ou permitir a reprodução do trabalho, mas que,
por sua extensão ou natureza, não seja adequado ao corpo principal do texto.

\begin{enumerate}
	\item Existe material produzido pelo autor que complementa o trabalho?
	\item Esse material é excessivamente extenso para permanecer no corpo do texto?
	\item Ele é importante para a reprodução do trabalho?
	\item Há códigos-fonte relevantes?
	\item Há tabelas extensas?
	\item Há deduções matemáticas detalhadas?
	\item Há questionários ou instrumentos de coleta desenvolvidos pelo autor?
	\item Há parâmetros completos de simulação?
	\item Há procedimentos experimentais detalhados?
	\item O leitor pode consultar esse material para aprofundar ou reproduzir o trabalho?
\end{enumerate}

{Exemplos de conteúdo de apêndices:}

\begin{itemize}
	\item Código desenvolvido;
	\item Dedução matemática extensa;
	\item Tabelas completas;
	\item Parâmetros de simulação;
	\item Questionários;
	\item \textit{Scripts};
	\item Procedimentos experimentais.
\end{itemize}